% math4cs-hw7-paths-cycles.tex

% !TEX program = xelatex
%%%%%%%%%%%%%%%%%%%%
% see http://mirrors.concertpass.com/tex-archive/macros/latex/contrib/tufte-latex/sample-handout.pdf
% for how to use tufte-handout
\documentclass[a4paper, justified]{tufte-handout}

\input{hw-preamble} % feel free to modify this file if you understand LaTeX well
%%%%%%%%%%%%%%%%%%%%
\title{计算机数学 (8-paths-cycles: 路径与圈)}
\me{魏恒峰}{hfwei@hnu.edu.cn}{}{}
\date{2026年5月15日}
%%%%%%%%%%%%%%%%%%%%
\begin{document}
\maketitle
%%%%%%%%%%%%%%%%%%%%
\noplagiarism % PLEASE DON'T DELETE THIS LINE!
%%%%%%%%%%%%%%%%%%%%
\begin{abstract}
\end{abstract}
%%%%%%%%%%%%%%%%%%%%
\beginrequired
%%%%%%%%%%%%%%%

%%%%%%%%%%%%%%%
\begin{problem}
  设 $G = (V, E)$ 是无向图 (不一定是简单无向图),
  其中 $|E| = m$。
  请证明\footnote{这也说明了, $G$ 中度数为奇数的顶点数目为偶数。},
  \[
    \sum_{v \in V} \deg(v) = 2m.
  \]
\end{problem}

\begin{proof}
\end{proof}
%%%%%%%%%%%%%%%

%%%%%%%%%%%%%%%
\begin{problem}
  \begin{enumerate}[(1)]
    \item 请证明: 每个 $u\-v$ 闭道路 (closed walk) 都包含一个 $u\-v$ 路径 (path)。
    \item 请证明: 每个长度为奇数的闭道路 (closed walk) 都包含一个长度为奇数的圈 (cycle)~\footnote{
    ``长度''就是所含边的条数。}。
  \end{enumerate}

  \vspace{1em}
  \noindent (提示: 可用数学归纳法。如果你使用数学归纳法, 请注意数学归纳法的书写规范。)
\end{problem}

\begin{proof}
\end{proof}
%%%%%%%%%%%%%%%

%%%%%%%%%%%%%%%
\begin{problem}
  设 $G$ 是一个简单无向图 (undirected simple graph) 且满足
  \[
    \delta(G) \ge k,
  \]
  其中 $k \in \mathbb{N}^{+}$ 为常数。
  请证明:
  \begin{enumerate}[(1)]
    \item $G$ 包含长度 $\ge k$ 的路径;
    \item 如果 $k \ge 2$, 则 $G$ 包含长度 $\ge k + 1$ 的圈。
  \end{enumerate}
  提示: By Extermality。
\end{problem}

\begin{proof}
\end{proof}
%%%%%%%%%%%%%%%

%%%%%%%%%%%%%%%
\begin{problem}
  考虑下图, 记为 $G$。
  \fig{width = 0.40\textwidth}{figs/Euler-graph}
  \begin{enumerate}[(1)]
    \item $G$ 是否是欧拉图? 请说明理由。
    \item 如果是欧拉图, 请将其分解为若干圈的组合, 并给出一个欧拉回路。
    　(可以自行为顶点编号, 也可以使用图上边的编号描述回路。)
  \end{enumerate}
\end{problem}

\begin{proof}
\end{proof}
%%%%%%%%%%%%%%%

%%%%%%%%%%%%%%%
\begin{problem}
  请判断下图 $G$ 是否是 Hamilton 图 (Hamiltonian graph), 并说明理由。
  \fig{width = 0.30\textwidth}{figs/Hamiltonian}
\end{problem}

\begin{proof}
\end{proof}
%%%%%%%%%%%%%%%

%%%%%%%%%%%%%%%
\begin{problem}
  在圆周上取 $n$ 个点。
  将每个点与圆周上两个方向上各自距离它最近的两个点相连,
  得到一个 $4$-正则图 $G_n$。

  \noindent 请证明: 当 $n \ge 5$ 时, $G_n$ 可以表示为两个 Hamilton 圈的并。
\end{problem}

\begin{proof}
\end{proof}
%%%%%%%%%%%%%%%

%%%%%%%%%%%%%%%%%%%%
% 如果没有需要订正的题目，可以把这部分删掉

\begincorrection
%%%%%%%%%%%%%%%%%%%%

%%%%%%%%%%%%%%%%%%%%
% 如果没有反馈，可以把这部分删掉
\beginfb

你可以写 (也可以发邮件)
\begin{itemize}
  \item 对课程及教师的建议与意见
  \item 教材中不理解的内容
  \item 希望深入了解的内容
  \item $\cdots$
\end{itemize}
%%%%%%%%%%%%%%%%%%%%
\end{document}